problem:  
Let \(M^4\) be a compact, simply connected Riemannian 4‑manifold with nonnegative sectional curvature and an effective isometric \(S^1\)-action. By the slice theorem, the local structure near each orbit is determined by the isotropy representation of \(S^1\) on the normal space. This local model defines, in particular, a 2‑dimensional quotient \(S^3/S^1_\text{iso}\), where \(S^1_\text{iso}\) is the isotropy subgroup, which serves as the space of directions in the orbit space \(X=M/S^1\).

**Open Rigidity Problem:**  
Does nonnegative sectional curvature restrict the *metric* of these local quotient spaces \(S^3/S^1_\text{iso}\) beyond what is implied by linear representation theory?  
Equivalently, suppose one fixes a smooth effective circle representation on \(\mathbb{R}^4\) with the induced linear quotient \(S^3/S^1_\text{iso}\). Among all nonnegatively curved metrics on \(M^4\) compatible with that orbit type, must the induced metric on \(S^3/S^1_\text{iso}\) (the link metric in \(X\)) be isometric to a round or canonical orbifold metric, or can curvature deformations alter the Alexandrov geometry of the orbit space without changing isotropy type?  

In other words: *Is the metric on the space of directions at singular orbits rigidly determined by the linear slice model under the assumption of nonnegative curvature, or can global curvature deformation produce new local Alexandrov metrics near the singular locus?*  

Why is it a "good" problem:  
This refined formulation reconciles the slice theorem with curvature constraints. It identifies a subtle, unexplored frontier: whether nonnegative sectional curvature forces local metric rigidity of the orbit space even when the isotropy representation (and thus topological type) is fixed. It directly connects Riemannian geometry (curvature bounds), transformation group theory (slice models and isotropy), and Alexandrov geometry (spaces of directions). Resolving it would clarify how global curvature interacts with local group symmetries—whether geometry around singular orbits is completely dictated by representation theory, or if curvature bounds allow previously unrecognized deformations of the local quotient metric. It is simple to state, lies at the core of the metric classification of nonnegatively curved manifolds, and its resolution would sharpen or overturn long‑standing assumptions about local rigidity under symmetry.